\cgeentry{cge_loop}{Solve the model for a list or grid of scenarios}

\cgesec{Description}

\cmd{cge_loop} (from version 0.9.9) solves the model for many scenarios in one
command, as a GAMS \cmd{loop} over a set of scenarios, or the GAMS scenario
solver GUSS, does. Each scenario sets the values of a few targets: elements of
parameters that \manref{cge_shock} changes, of user parameters, or of
variables fixed under the closure. It takes the shocks in place, starts from
the solution of the scenario before, and is stored as a simulation; report
measures and expressions are collected for every scenario, and the results can
be returned as a dataset. \manref{cge_gams} writes a loop as a GAMS program
that solves the scenarios in the same order.

\cgesec{Quick start}

A sweep over the tariff on manufactures, with government revenue and welfare
\begin{stcode}
cge_loop, name(tariff) values(tm(com_man) = 0(0.05)0.3) ///
    measures(YG, welfare equivalent_variation hhd)
\end{stcode}
Every combination of two tariffs, with the imports as a dataset
\begin{stcode}
cge_loop, name(two) values(tm(com_man) = 0 0.1 0.2; tm(com_agr) = 0 0.1) grid ///
    keep(QM) clear
\end{stcode}
Scenarios from a dataset with the variables \cmd{scenario}, \cmd{target} and
\cmd{value}
\begin{stcode}
cge_loop using scenarios.dta, name(policy)
\end{stcode}
The loop in GAMS, solved and compared with Stata
\begin{stcode}
cge_gams using gams_tariff, loop(tariff) replace run compare
\end{stcode}

\cgesec{Syntax}

\begin{stsyntax}
\cmdb{cge_loop} \texttt{,} \optarg{}{name}{name} \optarg{val}{ues}{spec} [\argu{options}]\\
\cmdb{cge_loop} \cmd{using} \argu{filename} \texttt{,} \optarg{}{name}{name} [\argu{options}]
\end{stsyntax}
\argu{spec} is \argu{target} \cmd{=} \argu{values}, separated by semicolons,
where \argu{values} is a numlist.

\begin{opttable}
\optgroup{Scenarios}
\optarg{}{name}{name} & the stub of the simulations, \argu{name}\cmd{1}, \argu{name}\cmd{2}, \dots\\
\optarg{val}{ues}{spec} & the values of each target\\
\opt{}{grid} & every combination of the lists of values\\
\opt{sc}{ale} & values as multiples of the targets' values in place\\
\optgroup{Solution}
\opt{}{cold} & solve each scenario from the base\\
\optarg{tol}{erance}{\#} & convergence tolerance of each solve; default \texttt{1e-10}\\
\optarg{iter}{ate}{\#} & Newton iterations of each solve; default 100\\
\optarg{par}{allel}{\#} & solve the scenarios in \# Stata processes (from version 0.19.0)\\
\optarg{stata}{exe}{path} & the Stata executable of the processes\\
\optarg{}{first}{\#} & the number of the first scenario; default 1\\
\optgroup{Results}
\optarg{meas}{ures}{list} & report measures or expressions, separated by commas\\
\optarg{}{keep}{names} & the variables of the dataset of results; all by default\\
\opt{}{clear} & replace the data in memory with the results\\
\opt{}{nostore} & do not store the simulations\\
\end{opttable}

\cgesec{Options}

\begin{optdesc}
\item[\optarg{}{name}{name}] names the loop; its scenarios are stored as the
simulations \argu{name}\cmd{1}, \argu{name}\cmd{2}, and so on, replacing
simulations of those names.

\item[\optarg{val}{ues}{spec}] gives each target its list of values, as
\cmd{values(tm(com_man) = 0 0.05 0.1; tm(com_agr) = 0 0.1 0.2)}. A target is
one element of a parameter that \manref{cge_shock} changes, such as
\cmd{tm(com_man)} or \cmd{trgr}, of a user parameter (\manref{cge_parameter}),
or of a variable fixed under the closure, such as \cmd{QFS(lab)}: the
instruments of \manref{cge_optimize}. By default the lists are taken in order,
one value of each per scenario, and must have the same length; a list of one
value applies to every scenario.

\item[\opt{}{grid}] makes every combination of the lists a scenario, the first
target varying slowest, as nested GAMS loops do.

\item[\opt{sc}{ale}] reads the values as multiples of the targets' values in
place.

\item[\cmd{using} \argu{filename}] reads the scenarios from a dataset with one
observation per scenario and target: \cmd{scenario}, a string or a number,
\cmd{target}, a string such as \cmd{tm(com_man)}, and \cmd{value}. Scenarios
and targets are taken in the order they first appear; a target that a
scenario leaves out keeps its value in place. The data in memory are left as
they are.

\item[\opt{}{cold}] solves each scenario from the base, with any starting
levels of \manref{cge_bound}, as separate solves would. By default each starts
from the solution of the scenario before.

\item[\optarg{meas}{ures}{list}] collects, for each scenario, report measures,
as \argu{table measure} [\argu{element}], or expressions of the user language
(\manref{cge_equation}), as in \manref{cge_optimize}.

\item[\opt{}{clear}] replaces the data in memory with one observation per
scenario and element of the variables kept, with \cmd{s}, \cmd{scenario},
\cmd{sim}, \cmd{converged}, the targets \cmd{t1}, \cmd{t2}, \dots, the measures
\cmd{m1}, \cmd{m2}, \dots, \cmd{variable}, \cmd{element1}, \cmd{element2},
\cmd{base} and \cmd{value}. \optarg{}{keep}{names} limits it to some
variables.

\item[\opt{}{nostore}] keeps the results without storing the simulations.

\item[\optarg{par}{allel}{\#}] (from version 0.19.0) solves the scenarios in
\# Stata processes at once, as GAMS's grid computing does: the model is saved,
the scenarios are split into contiguous groups, and each group is solved by
Stata in batch mode, which loads the model and runs \cmd{cge_loop} on it,
each scenario from the solution of the one before within its group; their
simulations are merged into the model in memory, with the same names and
results as one run. The processes run in the folder of Stata's temporary
files, whose files are removed afterwards. In batch mode on Windows, where
Stata ignores \cmd{winexec}, Stata 16 and later start the processes through
Python's \cmd{subprocess} module (from version 0.24.0; Stata's \cmd{python}
command, which needs a Python installation). This works only on some
versions of Stata: Stata 13 to 15 cannot start other processes from batch
mode, nor can Stata 16 or later without Python, and the scenarios are then
solved in the session, one after another, with a note that names the
version. Interactive Stata of any version starts the processes. Each process
starts Stata and loads the model, some seconds (about 14 on the machine of
the tests), so that parallel runs pay when the scenarios take longer than
that.

\item[\optarg{stata}{exe}{path}] gives the Stata executable of the parallel
processes; by default, that of the running edition in Stata's folder.

\item[\optarg{}{first}{\#}] (from version 0.19.0) numbers the scenarios from
\#, so that the simulations are \argu{name}\#, \dots; the parallel processes
use it.
\end{optdesc}

\cgesec{Remarks and examples}

A loop is a sequence of solves of the same model with different values of a
few parameters. Each scenario starts from the shocks in place: shocks applied
with \manref{cge_shock} before the loop are part of every scenario, and are
left as they were after it. Starting each scenario from the solution of the
one before makes a sweep over a parameter a path of small steps, as GUSS does,
so the loop takes fewer iterations than separate solves; a scenario that does
not solve from there is tried from the base and then by continuation, and one
that still fails is recorded, and the loop goes on.

The SAM is mapped and calibrated:
\stlog{loop_setup}

\cgesubsec{Example 1: a sweep over a tariff}

The tariff on manufactures from zero to 30 per cent, with government revenue,
the welfare of the household and spending on government purchases, an
expression:
\stlog{loop_sweep}
Each scenario is the simulation that \manref{cge_shock} and \manref{cge_solve}
would give on their own; \manref{cge_report} reports them as any other
simulations.

\cgesubsec{Example 2: a grid and the dataset of results}

Every combination of two tariffs, with the imports of each commodity as a
dataset:
\stlog{loop_grid}

\cgesubsec{Example 3: scenarios from a dataset}

A dataset gives the scenarios, one observation per scenario and target; here
as multiples of the values in place:
\stlog{loop_using}

\cgesubsec{Example 4: the loop in GAMS}

\manref{cge_gams}, \cmd{loop()}, writes the scenarios as the periods of a
dynamic run are written: the changes of each from the one before, solved in
order from the solution before, so that GAMS follows the same path:
\stlog{loop_gams}

\cgesec{Stored results}

\cmd{cge_loop} stores the following in \cmd{r()}:
\begin{storedres}
\resgroup{Scalars}
\res{r(N_scenarios)} & number of scenarios\\
\res{r(parallel)} & Stata processes that solved them (from version 0.19.0)\\
\res{r(N_converged)} & number of scenarios solved\\
\res{r(N_failed)} & number of scenarios not solved\\
\res{r(iterations)} & Newton iterations in all\\
\resgroup{Macros}
\res{r(name)} & the loop\\
\res{r(targets)} & its targets\\
\res{r(sims)} & its simulations, unless \opt{}{nostore}\\
\resgroup{Matrices}
\res{r(values)} & the targets' values, one row per scenario\\
\res{r(converged)} & 1 for each scenario solved\\
\res{r(measures)} & the measures, one row per scenario\\
\end{storedres}

\cgesec{Methods and formulas}

With targets $\theta_1, \dots, \theta_T$, their values in place $\theta^c$ and
the scenarios $s = 1, \dots, S$ with values $v_{s}$ (missing for a target that a
scenario leaves out, which keeps $\theta^c$), scenario $s$ solves the model of
\manref{cge_solve} with $\theta = v_s$ (or $\theta^c \odot v_s$ with
\opt{sc}{ale}) and the other shocks in place, from the point $x_{s-1}$, the
solution of the scenario before ($x_0$ the base, with the starting levels of
\manref{cge_bound}), or from the base with \opt{}{cold}. With \opt{}{grid},
the scenarios are the elements of $V_1 \times \dots \times V_T$ in
lexicographic order. A scenario that fails from $x_{s-1}$ is solved from the
base, then by continuation in 4 and then 16 steps (\manref{cge_solve}); the
next scenario starts from the last solution found.

\cgesec{Also see}

\begin{alsosee}
\seeentry{cge_shock}{Change a parameter or a fixed variable}
\seeentry{cge_solve}{Solve the model}
\seeentry{cge_optimize}{Choose policy instruments to maximise an outcome or reach targets}
\seeentry{cge_sensitivity}{Sensitivity analysis over elasticities}
\seeentry{cge_gams}{Write the model as a GAMS program, and solve it in GAMS}
\end{alsosee}
